\section{Boundaries and Future Routes}

The grammar leaves several routes outside the current claim boundary.  The signed-reversal program is now a public theorem route owned by P15 and remains pointer-only here.  P18 and P17 are future companion pointers, not admitted X11/X12 rows.  The exceptional \(S_6\) routes are closed in the bounded P07 and P33 companion packages but their theorems are not imported.  Higher reflection-shadow routes beyond the present exhibits and the blocked numerical bridge remain future work unless their own source locks, proof boundaries, and replay packages close.

The downstream P28 length-coloring theorem, P29 golden-ring refinement, P30 exceptional census, and P32 \(E_7\) placement study are likewise companion routes\cite{BabanskyyTypeBLengthColored2026,BabanskyyGoldenChannel2026,BabanskyyExceptionalFingerprints2026,BabanskyyGassmannE7Placements2026}.  They refine existing source rows or expose limitations of their descriptive fingerprints; none creates a new X-ID or transfers its theorem into P13.

\begin{table}[!htbp]
\centering\caption{Future routes and open debts kept outside the present claim boundary.}
\label{tab:future-routes}
\begin{tabular}{@{}L{0.20\textwidth}L{0.31\textwidth}L{0.39\textwidth}@{}}
\toprule
route or debt & current role & admission condition before promotion \\
\midrule
signed-reversal route & public companion theorem route & P15 owns the theorem; kept pointer-only here \\
P18 nonconvex route & future companion pointer & remains \texttt{assigned\_p13\_label=null}; a separate P13 admission gate would be required \\
P17 motion/contact route & future companion pointer & finite CL3 candidate does not itself create an admitted P13 row \\
higher reflection-shadow routes beyond the bounded exhibits & future grammar branches & need source locks and non-overlap with this paper's dihedral, \(G_2\), projective \(F_4\), projective \(H_3\), projective \(H_4\), projective \(D_5\), projective \(E_6\), projective \(E_7\), and projective \(E_8\) exhibit claims \\
M19 and \texttt{lat565} & controlled X05 exhibit & \texttt{lat565} is local \(32+8+22\) only; broader out-of-\(G_0\) mechanisms remain outside this paper \\
shifted-diagonal row & P12 control family (pointer; not a P13 exhibit) & all-\(n\) statement is an open obligation, homed at EXT-NEG-19 \\
5568 bridge candidate & blocked numerical coincidence/bridge candidate & requires an identified morphism or explicit non-bridge closeout \\
exceptional \(S_6\) benchmark & public companion certificate routes & P07/P33 own their bounded claims; no theorem import into P13 \\
source-quotient coupling route & closed guardrail exhibit (\cref{prop:source-coupling}) & any transfer-law promotion needs genuinely non-split source families and its own admission gate; the mechanisms are classical \\
\bottomrule
\end{tabular}\end{table}

The modular centered row is a compact channel exhibit only.  It records that rational centered rank is not enough once finite-field reduction is admitted: the grammar must name the characteristic, the augmentation channel, the reduced quotient when p divides the carrier size, and the lattice-qualified Smith form.  It does not claim an all-layer Type-A modular theorem, a Latin/Sudoku/M19 mechanism, a modular Specht novelty theorem, a classifier, or a tiling result.

The X05/M19 exhibit has its own boundary.  This paper may use it to show that \(\PhiFCIG\) can recover a recurring hidden group in a centered-operator rank locus and can record, for \texttt{lat565} only, a finite \(32+8+22\) centered-certificate stratification.  This paper does not claim a general Sudoku or Latin theorem, a classification of X05 rank drops or out-of-\(G_0\) objects, a second hidden group for the full \texttt{lat565} rank locus, or any tiling consequence.  The recurring hidden group recovered here is the affine symmetry group \(3^2{:}D_4\) of the natural \(3\times 3\) (Lo Shu) value grid, certified single-base in the companion Sudoku package (\(K_L\neq K_R\)); its recurrence as a single universal \(G_0\) across the census is recorded here as computational evidence, not a certified cross-grid identity.

The dihedral row also has its own boundary.  It is a rank-two vertex-shadow exhibit inside \(S_m\), not a general Coxeter or Weyl theory.  The special case \(I_2(6)\) is not used as a theorem about \(G_2\), and the lattice-qualified Smith data in \cref{tab:dihedral-witnesses,tab:phi-payoff} is not promoted to a classifier.

The length-colored \(G_2\) row is likewise bounded.  It uses the twelve-root finite shadow only with its source partition \(\Omega_{\mathrm{s}}\sqcup\Omega_{\mathrm{l}}\).  The projective quotient and the uncolored dodecagon are controls, not replacement sources, because both forget the short/long channel.  This paper does not claim a general Weyl-group grammar, a projective theorem, or any tiling consequence from the \(G_2\) exhibit.

The projective \(F_4\) row is bounded in the same way.  It uses only the 24-point projective root shadow \(\Pi_{24}\), the long/short source partition, the projective Weyl layer \(G_{\mathrm{proj}}\), and the source matroid automorphism layer \(G_{\mathrm{mat}}\).  It does not classify \(F_4\) rows, block/triality rows, flat-incidence rows, 24-cell geometry, Weyl groups, or tilers, and it does not claim the \(F_4\) matroid automorphism group as a new discovery.

The projective \(H_3\) row is bounded in parallel.  It uses only the 15-point projective root shadow \(\Pi_{15}\), the flat-size edge-coloring channel, the Gram-angle edge-coloring channel, and the ordinary \(S_{15}\) Specht scout as descriptive fingerprint data.  It does not prove a general \(H_3\), non-crystallographic Coxeter, root-system matroid, classifier, or tiling theorem.  The public source posture is the bundled finite coordinate/source-channel replay over \(\mathbb Q(\phi)\); no external A--O notation is used as a page-locked proof source here.


The projective \(H_4\) row is even more tightly bounded.  It uses only the 60 antipodal pairs in the \(H_4\) root system, viewed as a projective 600-cell root shadow, and only the two pair-color channels recorded in the selected source package: rank-two flat size and squared Gram value.  It does not claim a new computation of \(\operatorname{Aut}(M(H_4))\), a full \(H_4\) theory, a 600-cell classification, an E-type or icosian theorem, a Weyl/Coxeter theorem, a classifier, or a tiling result.  The public posture is the bundled de-labelled H4 replay/package, not a general H4 engine.


The projective \(D_5\) row is a control exhibit only.  It uses the 20 antipodal projective pairs in the \(D_5\) root system, the pair-level Gram/flat-size channel, the coordinate \(D_4\) negative-control family, and the 40 projective \(A_2\) 3-flat incidences.  It does not claim a standalone \(D_5\) theorem, a general Type-E theorem, an \(E_6/E_7/E_8\) transfer, a matroid automorphism classification, a classifier, or a tiling result.  The \(A_2\) incidence is a triple/flat-incidence channel, not another pair channel.

The projective \(E_6\) row is also bounded.  It uses only the 36 antipodal projective pairs in the \(E_6\) root system, the projective root graph, the flat-size and circuit channels, and the source summands \(U_{20}\) and \(U_{15}\).  It does not claim a standalone \(E_6\) theorem, a Type-E family theorem, a Schlaefli/double-six/27-line/cubic-surface route, a new root-system matroid result, a classifier, or a tiling consequence.  The parabolic witness is explicitly not a same-type or same-orbit indistinguishability witness.
The projective \(E_7\) row is bounded in the same way.  It uses only the 63 antipodal projective pairs in the \(E_7\) root system, the Gram-square/projective root graph, the rank-two flat-size pair channel, the source summands \(U_{27}\) and \(U_{35}\), and selected embedded simple-parabolic rows.  It does not claim a standalone \(E_7\) theorem, a theorem about \(\operatorname{Sp}_6(2)\), a reflection-subgroup classification, a root-system matroid novelty result, an \(E_8\) or Type-E family theorem, a 56-weight/Gosset/minuscule route, a classifier, or a tiling consequence.  The parabolic separator is explicitly an embedded-row fingerprint witness, not a classification of parabolics.
The projective \(E_8\) row is bounded likewise.  It uses only the 120 antipodal projective pairs in the \(E_8\) root system, the Gram-square/projective root graph, the rank-two flat-size pair channel, the source summands \(U_{35}\) and \(U_{84}\), the standard simple-parabolic fusion replay, and selected coordinate-\(D_8\) subsystem rows.  It does not claim a standalone \(E_8\) theorem, a Type-E family theorem, a full reflection-subgroup or root-subsystem classification, a new Dynkin--Borel--de Siebenthal--Oshima result, a root-system matroid novelty result, a classifier, or a tiling consequence.  The subsystem separator is explicitly a bounded source-row fingerprint witness, not a classification of \(E_8\) subsystems.
The shifted-diagonal material and the 5568 bridge candidate remain fenced.  The shifted-diagonal row is a pointer to the P12 control family, not a P13 exhibit; its all-\(n\) statement is an open proof obligation, consolidated at the EXT-NEG-19 home.  The numerical bridge involving 5568 is not used as a theorem or as evidence of an identified morphism.

This boundary is part of the result of the interface.  The grammar is useful only if it admits source-locked finite rows while refusing rows whose finite shadow, representation channel, or fingerprint package is not yet defined.  If later drafting removes the worked examples in \cref{tab:type-a-rankmass,tab:type-a-value-witness,tab:modular-channel-witnesses,tab:signed-rankmass,tab:signed-value-witness,tab:dihedral-witnesses,tab:g2-source-channel,tab:f4-source-layer,tab:h3-source-channel,tab:h4-source-channel,tab:d5-root-shadow-control,tab:e6-root-shadow-control,tab:e7-root-shadow-control,tab:e8-root-shadow-subsystem-separator,tab:phi-payoff,tab:m19-census,tab:abd-contrast,tab:future-routes}, the honest fallback is a shorter interface note rather than a full paper.